\section{Weil classes and their detecting cycles}
\label{sec:weil}
This section studies the quantitative form of algebraicity for the two-dimensional space of Weil classes. Let $K$ be an imaginary quadratic field acting on a polarized abelian variety $A$ of dimension $2n$, with signature $(n,n)$ and with the Rosati involution restricting to complex conjugation on $K$. Its Weil space is the rational plane
\[
 W_K(A)=\bigwedge_K^{2n}H^1(A,\Q)\subset H^{2n}(A,\Q).
\]
It consists of Hodge classes. Write $e_W$ for its projector in the eigenspace decomposition for the $K$-action, and write $\delta_W(A,\theta)$ for the least degree of an irreducible reduced $n$-dimensional subvariety $Y$ with $e_W\cl(Y)\ne0$, or $+\infty$ if no such subvariety exists. Such a subvariety will be called a detector.

The principal open question for the Weil-type part of this article is whether $\delta_W$ is bounded on the CM points in every connected split Weil family, in every dimension and over every imaginary quadratic field. Here ``split'' means that the rational Hermitian space of dimension $2n$ and signature $(n,n)$ is a sum of hyperbolic planes. Its determinant is $(-1)^n$ modulo norms from $K$; the positive signed discriminant is therefore $1$ modulo norms. We use ordinary intersection degrees throughout.

\subsection{One detector gives the entire plane}

\begin{theorem}[Quadratic extraction]\label{thm:extraction}
Let $n\geq1$ and let $A$, $\theta$ and $K$ be as above, with $\theta$ integral and ample. Let $a\in K$ act by an endomorphism of $A$, assume $a/\bar a$ is not a root of unity, and put $q=\Nm_{K/\Q}(a)$. For a nonzero $t\in W_K(A)$, choose $\delta\in\Z_{>0}$ with $\delta t$ integral and set
\[
 T=(-1)^n\int_A t^2>0.
\]
Then
\begin{equation}\label{eq:sharper-extraction}
 \delta_W\leq c_t\leq
 \left\lceil\frac{2\,16^n}{n!}\right\rceil
 \delta\lceil\sqrt T\rceil\,
 q^{n(2n+1)}\delta_W^2.
\end{equation}
The assertion is valid for every polarization type. In particular it implies
\begin{equation}\label{eq:original-extraction}
 c_t\leq
 \left\lceil\frac{8\,16^n}{n!}\right\rceil
 \delta\lceil\sqrt T\rceil\,
 q^{2n(n+1)}\delta_W^2.
\end{equation}
If no detector exists, both $\delta_W$ and $c_t$ are infinite. Thus the existence of one detector is equivalent to algebraicity of the entire rational Weil plane.
\end{theorem}

The proof has two parts. A polynomial in endomorphism pullbacks removes every middle eigenspace except the Weil plane. Two integral vectors in that positive plane then span a sublattice of finite index. Clearing that index, rather than estimating an arbitrary inverse matrix, keeps the cost quadratic.

\begin{lemma}[Norm and degree]\label{lem:weil-norm}
The Weil plane is primitive for $\theta$, and
\[
 B(u,v)=(-1)^n\int_Au\smile v
\]
is a positive definite integral form on its integral lattice. With
$v_\theta=\int_A\theta^{2n}/(2n)!\in\Z_{>0}$ and
$\|u\|^2=B(u,u)$, every effective $n$-cycle $Y$ satisfies
\[
 \|e_W\cl(Y)\|\leq\frac{\deg_\theta Y}{n!\sqrt{v_\theta}}.
\]
For any chosen signed presentation $C=C^+-C^-$ whose class lies in $W$, its norm is at most
$(\deg C^++\deg C^-)/(n!\sqrt{v_\theta})$.
\end{lemma}

\begin{proof}
Over $\C$ write $V=H^1(A,\C)=V_\sigma\oplus V_{\bar\sigma}$ for the two eigenspaces of $K$, each of dimension $2n$. The signature hypothesis makes the extreme determinant lines
\[
 W_\C=\bigwedge^{2n}V_\sigma\oplus\bigwedge^{2n}V_{\bar\sigma}
\]
of Hodge type $(n,n)$. Rosati compatibility gives $a^*\theta=q\theta$, so $\theta$ lies in $V_\sigma\otimes V_{\bar\sigma}$. Its product with either extreme line vanishes for dimension reasons. This proves primitivity.

The primitive Hodge-star formula gives $*w=(-1)^nw$ for $w\in W_\R$, and consequently $B(w,w)=\int_A w\wedge*w>0$. This is the middle-degree Hodge--Riemann form, with the stated sign \cite[Theorem~5.42 and (5.43)]{Ballmann}. The cup product is integral on integral cohomology.

Use the translation-invariant K\"ahler metric determined by $\theta$. Harmonic forms are translation-invariant, so a vector $w$ of cohomological norm one has pointwise norm $v_\theta^{-1/2}$. Wirtinger's formula gives $\operatorname{vol}(Y)=\deg_\theta(Y)/n!$, including for singular $Y$ by integration on its regular locus. Therefore
\[
 \left|\int_Yw\right|\leq\frac{\deg_\theta Y}{n!\sqrt{v_\theta}}.
\]
The cup-product projection onto $W$ is the metric projection: for $w\in W$, its metric pairing with any middle-degree class $u$ is $\int_Au\wedge*w=B(u,w)$. Taking the supremum over unit $w$ proves the assertion for an effective cycle, and the triangle inequality proves it for a signed cycle.
\end{proof}

\begin{proof}[Proof of Theorem~\ref{thm:extraction}]
The trace of $a$ is integral, and hence $\bar a=\operatorname{Tr}_{K/\Q}(a)-a$ is also an endomorphism. On $H^{2n}$ set
\[
 S=a^*+\bar a^*,\qquad
 \beta_j=q^j\bigl(a^{2n-2j}+\bar a^{2n-2j}\bigr)
 \quad(0\leq j\leq n).
\]
Here $S$ denotes a sum of pullback operators. On the summand with $j$ factors from $V_{\bar\sigma}$ and $2n-j$ from $V_\sigma$, its eigenvalue is $\beta_j=\beta_{2n-j}$. These numbers are integers, and $|\beta_j|\leq2q^n$. If $z=a/\bar a$, equality of two distinct $\beta_j$, for $0\leq j\leq n$, would give $z^{j-k}=1$ or $z^{2n-j-k}=1$. This is excluded by the hypothesis. Thus
\[
 P(X)=\prod_{j=1}^n(X-\beta_j)\in\Z[X],\qquad
 P(S)=P(\beta_0)e_W,
\]
and $P(\beta_0)\ne0$. The middle eigenspaces are orthogonal to the extreme lines because a nonzero top wedge must have exactly $2n$ factors from each $K$-eigenspace. This verifies that this operator is the projector specified above, multiplied by a nonzero integer.

For a chosen signed cycle presentation put
\[
 \mathcal D(C)=\deg_\theta C^++\deg_\theta C^-.
\]
Pullback by $a$ preserves effectivity and multiplies the degree of an $n$-cycle by $q^n$: indeed $\deg(a)=q^{2n}$ and $a^*\theta=q\theta$. It follows that
\[
 \mathcal D((S-\beta_j)C)
 \leq(2q^n+|\beta_j|)\mathcal D(C)
 \leq4q^n\mathcal D(C).
\]
Choose an irreducible detector $Y$ of degree $d=\delta_W$. The integral signed cycle $C=P(S)Y$ has nonzero class $u\in W$ and has a presentation with
\[
 \mathcal D(C)\leq L:=4^nq^{n^2}d,
 \qquad \|u\|\leq L/n!.
\]
Its pullback $a^*C$ has mass at most $q^nL$ and class $v=a^*u$ with $\|v\|=q^n\|u\|$. The vectors $u,v$ are independent over $\Q$: the eigenvalues $a^{2n},\bar a^{2n}$ of $a^*|_W$ are nonreal conjugates, so this rational operator has no rational eigenline.

Let $\mathcal W=W\cap H^{2n}(A,\Z)$. It is a saturated rank-two lattice with integral positive Gram matrix. Its Euclidean covolume $c$ satisfies $c^2\in\Z_{>0}$, so $c\geq1$. Give $W_\R$ the orientation for which $\det(u,v)>0$, and let
\[
 m=[\mathcal W:\Z u+\Z v].
\]
Then $mc=\det(u,v)\leq\|u\|\|v\|$, and $m\mathcal W\subset\Z u+\Z v$. Therefore there exist integers $r,s$ such that
\[
 m\delta t=ru+sv.
\]
Taking determinants with $v$ and with $u$ gives
\[
 r=\frac{\delta\det(t,v)}c,\qquad
 s=\frac{\delta\det(u,t)}c,
\]
and hence
\[
 |r|\leq\delta\sqrt T\,\|v\|,
 \qquad |s|\leq\delta\sqrt T\,\|u\|.
\]
The signed cycle $rC+s a^*C$ represents $M t$, where $M=\delta m$ is a positive integer. Grouping positive and negative coefficients gives effective integral cycles satisfying
\[
 M\leq\frac{\delta q^nL^2}{(n!)^2},\qquad
 \deg Z^\pm\leq|r|L+|s|q^nL
 \leq\frac{2\delta\sqrt T\,q^nL^2}{n!}.
\]
Substitution of $L$ proves \eqref{eq:sharper-extraction}. Since $q\geq1$, it also proves \eqref{eq:original-extraction}.

Finally, if $Mt=\cl(Z^+)-\cl(Z^-)$, at least one irreducible component of one of these cycles has nonzero Weil projection. Its degree is no greater than that of the corresponding effective cycle. Thus $\delta_W\leq c_t$. If no detector exists, this observation excludes every finite representative.
\end{proof}

\begin{remark}[An endomorphism available in every order]
If $K=\Q(\sqrt{-m})$, with $m>0$ squarefree, and the endomorphism order contains $\Z+f\mathcal O_K$, one may take $a=f(2+\sqrt{-m})$. Then $q=f^2(m+4)$. Outside $\Q(i)$ and $\Q(\sqrt{-3})$, the only roots of unity in $K$ are $\pm1$, whereas $a/\bar a$ is neither. In the two exceptional fields its real part is respectively $3/5$ and $1/7$, so it is again not a root of unity.
\end{remark}

\subsection{Bounded detector loci}

\begin{proposition}\label{prop:detector-loci}
In a connected polarized family with a fixed integral quadratic action of signature $(n,n)$, the locus $\{s:\delta_W(s)\leq D\}$ is Zariski closed. The following are equivalent:
\begin{enumerate}[label=\textnormal{(\roman*)}]
\item every Weil class on every fibre is algebraic;
\item $\delta_W$ is bounded on every fibre;
\item $\delta_W$ is bounded on one Zariski-dense set.
\end{enumerate}
For a special Weil family these conditions are also equivalent to one uniform detector-degree bound at infinitely many good reductions of every CM fibre, allowing the auxiliary prime and geometric place to vary.
\end{proposition}

\begin{proof}
On the projective parameter space of effective cycles of degree at most $D$, the projected universal class $e_W\gamma$ is a section of the pulled-back Weil local system. Its nonvanishing locus is a union of connected components, hence is also closed and proper over the base. Its image is exactly the stated locus: if an effective sum detects, some irreducible component detects. This proves closedness.

If every Weil class is algebraic, every fibre has a detector. The closed loci therefore cover the complex base; Baire's theorem gives one bound. Conversely, a bound on a dense set gives a bound everywhere by closedness, and Theorem~\ref{thm:extraction} makes every Weil class algebraic. Alternatively, the integral Weil local system has finite monodromy, because it preserves the positive integral form $B$; after a finite \'etale cover it has a basis of global rational sections to which Theorem~\ref{thm:family} applies.

For the arithmetic assertion, extend the quadratic endomorphisms over an arithmetic open of each CM fibre. The projector is a rational polynomial in their pullbacks and commutes with specialization. The proof of Proposition~\ref{prop:lift} therefore lifts a bounded-degree special cycle with nonzero projected class to a bounded-degree characteristic-zero detector, at one place outside the finite exceptional set. Nonvanishing is preserved by the specialization isomorphism. In the reverse direction, descend and spread a characteristic-zero detector as in Proposition~\ref{prop:spread}. This gives detectors at all but finitely many good places.
\end{proof}

\subsection{An explicit Gaussian extraction constant}

\begin{proposition}\label{prop:gaussian-extraction}
On a Gaussian Weil fourfold with integral $\Z[i]$-action, a detector of degree $D$ gives
\[
 c_t\leq61\,440\,\delta\lceil\sqrt T\rceil D^2.
\]
\end{proposition}

\begin{proof}
Put $I=[i]^*$ and $J=[1+i]^*$. On the five summands of $H^4$ indexed by the number of factors from one $K$-eigenspace, direct substitution gives
\[
 (1+I)(4-J)=16e_W.
\]
For a detector $Y$, set $y=16e_W[Y]$. It is represented by
\[
 4Y+4[i]^*Y-[1+i]^*Y-[i(1+i)]^*Y,
\]
whose two effective parts each have degree $8D$. The norm lemma gives $d:=B(y,y)\leq64D^2$. The operator $A=[2+i]^*$ has eigenvalues $-7\pm24i$ on $W$. Hence $z=(A+7)y$ satisfies
\[
 B(y,z)=0,\qquad B(z,z)=576d,
\]
and has signed effective parts of degree at most $256D$. For the integers $b_1=B(\delta t,y)$ and $b_2=B(\delta t,z)$, orthogonal expansion gives
\[
 (576\delta d)t=576b_1y+b_2z.
\]
Write $t=x y/\sqrt d+x' z/(24\sqrt d)$, so $x^2+(x')^2=T$. The degree of each signed part is at most
\[
 \delta D\sqrt d\,(4608|x|+6144|x'|)
 \leq7680\delta\sqrt T D\sqrt d
 \leq61\,440\delta\sqrt T D^2.
\]
The multiplier is at most $36\,864\delta D^2$, which is bounded by the asserted expression.
\end{proof}

\section{Stabilization after Schoen}
\label{sec:stabilization}

This section explains how bounded detectors in the next split dimension control every determinant in the preceding dimension. Schoen's product argument is stated for fourfolds and sixfolds in \cite[Proposition~10, pp.~332--333]{Schoen98}; the general product and contraction statement is also recorded by Voisin \cite[Lemma~2.9]{VoisinB}. We include the degree estimate, for which the polarization and the chosen effective divisors matter.

\begin{theorem}[Qualitative stabilization]\label{thm:stabilization}
Fix an imaginary quadratic field $K$ and $n\geq1$. Suppose that the detector degree is bounded in every connected split Weil family of dimension $2n+2$. Then it is bounded in every connected Weil family of dimension $2n$, of every determinant. The bound may depend on the family and its polarization. No class-number hypothesis is required.
\end{theorem}

\begin{proof}
The determinant class of a fixed connected Weil family is constant. Write it as $(-1)^n d$ modulo norms, with $d=r/s>0$ and $r,s$ positive integers. On $B=E_K\times E_K$, where $E_K=\C/\mathcal O_K$, use the action $a\cdot(x,y)=(ax,\bar a y)$ and the polarization $\lambda=r\vartheta_1+s\vartheta_2$. Its Hermitian determinant is $-rs$ modulo norms, hence $-d$ because $rs=d s^2$. The product $A\times B$, with polarization $\Theta=\theta+\lambda$, has determinant $(-1)^{n+1}d^2$, so its Hermitian space is split by the rational classification \cite[Proposition~4.1 and Corollary~4.2]{Deligne82}. All this data is fixed in the family.

The extreme $K$-eigenlines give an inclusion
\[
 W_K(A\times B)\subset W_K(A)\otimes W_K(B).
\]
Choose an integral basis $\eta_1,\eta_2$ of $W_K(B)\cap H^2(B,\Z)$. The Lefschetz $(1,1)$-theorem makes them divisor classes. For a sufficiently large fixed $N$, the classes
\[
 d_0=N\lambda,\qquad d_1=N\lambda+\eta_1,
 \qquad d_2=N\lambda+\eta_2
\]
are very ample. Choose $b_B>0$ such that $b_B\lambda-d_j$ is nef for each $j$.

Let $Y$ be an effective detector on the split product, of degree at most $D$. Choose effective divisors $D_j$ in the classes $d_j$, translating them to meet $Y$ properly after pullback, and set
\[
 C_j=p_{A*}(Y\cdot p_B^*D_j).
\]
These are effective $n$-cycles. If all three had zero Weil projection, the projection formula and $d_i-d_0=\eta_i$ would make $[Y]$ orthogonal to every $w\otimes\eta_i$, with $w\in W_K(A)$. It would then be orthogonal to $W_K(A\times B)$, contradicting detection. Thus some $C_j$ detects. Mixed nef intersections give
\[
 \deg_\theta C_j
 \leq b_B\int_Yp_A^*\theta^n p_B^*\lambda
 \leq\frac{b_B}{n+1}\int_Y\Theta^{n+1}
 \leq\frac{b_BD}{n+1}.
\]
The factor $n+1$ is the coefficient of $\theta^n\lambda$ in the binomial expansion. Taking an irreducible detecting component proves a uniform bound on $A$.

If the hypothesis is formulated only for principal polarizations, pass to a fixed isogenous principally polarized split family. The choice of a maximal isotropic subgroup of the polarization kernel becomes constant after a finite \'etale cover. The rational $K$-action transports across the isogeny, and an order in $K$ acts integrally. For an isogeny $f:C\to C'$ with $f^*\theta'=k^2\Theta$, a middle-dimensional cycle satisfies
\[
 \deg_\Theta f^*Y'
 =\frac{\deg f}{k^{2n+2}}\deg_{\theta'}Y'.
\]
This is a constant factor in the family. Pullback preserves a nonzero Weil projection, so the preceding contraction applies. The finite cover is surjective, and therefore the bound descends as a fibrewise existence assertion.
\end{proof}

\begin{theorem}[An explicit contraction bound]\label{thm:explicit-stabilization}
Assume $\mathcal O_K=\Z[\tau]$ has class number one and acts integrally on a polarized Weil $2n$-fold $(A,L)$. Put $\Delta=\chi(L)$ and
$\kappa_K(\Delta)=\max\{1,\Nm(\tau)/\Delta\}$. If every split Weil $(2n+2)$-fold of polarization product $x$ has a detector of degree at most $D_{n+1}(x)$, then
\[
 \delta_W(A,L)\leq\frac{\kappa_K(\Delta)}{n+1}D_{n+1}(\Delta^2).
\]
For $K=\Q(i)$ or $\Q(\omega)$ choose $\tau=i$ or $\omega$, respectively; then $\kappa_K(\Delta)=1$. The premise ranges over all polarization types of the specified product, not only principal polarizations.
\end{theorem}

\begin{proof}
Class number one makes $H_1(A,\Z)$ free over $\mathcal O_K$. Normalize the Hermitian form associated to the alternating Riemann form $E$ by
\[
 h(x,y)=E(x,\tau y)-\bar\tau E(x,y).
\]
In an integral $\mathcal O_K$-basis it has determinant $(-1)^n\Delta$. Indeed for one diagonal block $h=(a)$, the alternating form on $(v,\tau v)$ has matrix $\left(\begin{smallmatrix}0&a\\-a&0\end{smallmatrix}\right)$. Both sides of $\det_\Z E=(\det_Kh)^2$ transform by the same norm factor under a change of $K$-basis. The sign is fixed by signature and $|\det_Kh|=\chi(L)$.

Take $B=E_K^2$ with balanced action and $\lambda=\vartheta_1+\Delta\vartheta_2$, where $\vartheta_1$ is the class of $\{0\}\times E_K$ and $\vartheta_2$ that of $E_K\times\{0\}$. Its determinant is $-\Delta$, so $A\times B$ is split and has polarization product $\Delta^2$. Let $Y$ be a detector there, of degree at most $D=D_{n+1}(\Delta^2)$ for $\Theta=\theta+\lambda$, and put
\[
 I=\int_Y\theta^n\lambda.
\]
All terms in the expansion of $\Theta^{n+1}$ have nonnegative intersection with $Y$, since both factors are nef. Hence $(n+1)I\le D$.

Intersect with general translates of the coordinate elliptic divisors and push to $A$, as in the preceding proof. Each contraction has degree at most $I$, since $\vartheta_j\le\lambda$ in the nef order. If either has nonzero Weil projection, an irreducible component gives the result.

Suppose both projected contractions vanish. For $u=1,\tau$ use the two graphs $\Gamma_u=\{x=u y\}$ and $\Gamma_{-u}$. Their classes satisfy
\begin{equation}\label{eq:opposite-graphs}
 [\Gamma_u]+[\Gamma_{-u}]
 =2(\vartheta_1+\Nm(u)\vartheta_2)
 \le2\kappa_K(\Delta)\lambda
\end{equation}
in the nef order. The equality follows either by adding the Hermitian matrices
$\left(\begin{smallmatrix}1&\mp u\\\mp\bar u&\Nm(u)\end{smallmatrix}\right)$ or from the equation $x\mp uy=0$. The axes and the graphs of $1,\tau$ span $\operatorname{NS}(B)_\Q$: the axes give the diagonal entries and $1,\tau$ give a rational basis for the off-diagonal entry. Their off-diagonal parts span $W_K(B)$.

Because $Y$ detects $W_K(A\times B)$, contraction against some element of $W_K(B)$ has nonzero projection to $W_K(A)$. Thus one of the two graph pairs has a detecting member. Its two projected contractions are opposite by \eqref{eq:opposite-graphs} and the vanishing of the axis contractions, so both detect. Choose the member of smaller degree. Its degree is at most the average, which by \eqref{eq:opposite-graphs} is at most $\kappa_K(\Delta)I$. One of its irreducible components detects and has no larger degree. This proves the bound.
\end{proof}

\begin{remark}[The polarization range]
Full integral maximal-order action and principal polarization over a class-number-one field force $\Delta=1$ and determinant $(-1)^n$, hence the split determinant class. Statements covering every determinant must therefore allow nonprincipal polarizations or rational actions transported through isogenies. The preceding theorem does so explicitly.
\end{remark}
